\section{Univalent maps and a dyadic transfer operator}
\label{sec:preliminaries}

We first place normalized univalent maps of the upper half-plane in a
compact space.  Rescaling the source about a point then defines a positive
operator weighted by the inverse square of the derivative.  The main
conclusion of this section is a uniform comparison between powers of that
operator and disk integral means.  Later sections will bound the operator's
growth rate.

\subsection{Classical estimates and the normalized compact class}

We use the classical coefficient, distortion, growth, and quarter-disk
estimates in the following precise forms; see Duren
\cite[Chapter~2, especially \S2.3]{Duren1983}. No boundary regularity
or boundedness of the image is required.

\Needspace{12\baselineskip}
\begin{lemma}[Classical univalence estimates]
\label{lem:classical-univalence}
If $f(z)=z+a_2z^2+\cdots$ belongs to $\SSS$, then $|a_2|\le2$ and,
for $|z|=r<1$,
\begin{equation}
 \frac{1-r}{(1+r)^3}\le |f'(z)|\le\frac{1+r}{(1-r)^3},
 \qquad \frac r4\le |f(z)|\le\frac{r}{(1-r)^2}.
 \label{prelim:disk-bounds}
\end{equation}
Every univalent holomorphic map $g$ on an open disk $B(z_0,R)$ satisfies
\begin{equation}
 B\bigl(g(z_0),R|g'(z_0)|/4\bigr)\subset g(B(z_0,R)).
 \label{prelim:quarter}
\end{equation}
\end{lemma}

\begin{proof}
The coefficient and derivative bounds are the classical estimates just
cited. The classical lower growth bound $r/(1+r)^2$ implies the weaker
bound $r/4$ recorded here. A univalent holomorphic map has nonzero
derivative, so for the last assertion the map
\[
 \zeta\longmapsto
 \frac{g(z_0+R\zeta)-g(z_0)}{R g'(z_0)}
\]
belongs to $\SSS$. Applying the Koebe quarter theorem to this normalized
map gives \eqref{prelim:quarter}.
\end{proof}

Write
\[
 \HH=\{x+iy:x\in\mathbb R,\ y>0\},\qquad
 m(z)=\frac{z-i}{z+i},
\]
and let
\[
 \KK=\{F:\HH\to\CC\text{ univalent}:F(i)=0,\ F'(i)=1\}.
\]
We give $\KK$ the topology of locally uniform convergence.  Since
$m'(i)=1/(2i)$, its correspondence with the disk class is exactly
\begin{equation}
 F(z)=2i f(m(z)),\qquad
 f(\zeta)=\frac{F(m^{-1}(\zeta))}{2i}.
 \label{prelim:normalization}
\end{equation}
For $F\in\KK$, define the holomorphic function
\[
 Q_F(z)=\frac1{F'(z)}.
\]

\begin{lemma}[Compactness and uniform local bounds]
\label{lem:compact-class}
The space $\KK$ is compact and metrizable.  For every integer $j\ge0$,
the evaluations $(F,z)\mapsto F^{(j)}(z)$ and
$(F,z)\mapsto Q_F^{(j)}(z)$ are jointly continuous on $\KK\times\HH$
and bounded on $\KK\times E$ for every compact $E\subset\HH$.
For every $C_0\ge0$,
\begin{equation}
 |Q_F(x+iy)|\le\frac{(C_0^2+4)^2}{y}
 \quad\bigl(F\in\KK,\ |x|\le C_0,\ 0<y\le1\bigr).
 \label{prelim:strip-bound}
\end{equation}
In particular the bound is $64/y$ when $|x|\le2$.  Also,
\begin{equation}
 |F(z)|\ge\frac1{16}
 \quad\text{if }F\in\KK\text{ and }|m(z)|\ge\frac12.
 \label{prelim:zero-separation}
\end{equation}
\end{lemma}

\begin{proof}
The upper growth bound and \eqref{prelim:normalization} give local uniform
boundedness.  Montel's theorem therefore gives a locally uniformly
convergent subsequence from every sequence in $\KK$; see
\cite[Sections~1.1 and~1.3]{Duren1983} for this normal-family argument and
the univalent-limit form of Hurwitz's theorem.  Cauchy's formula preserves
the derivative normalization, so a limit $F$ has $F'(i)=1$ and is
nonconstant.  For any fixed $z_0\in\HH$, apply Hurwitz's theorem to
$F_n(z)-F_n(z_0)$ on the connected domain $\HH\setminus\{z_0\}$.
The limit is not identically zero and has no zero there.  Thus $F$ is
injective and belongs to $\KK$.  The compact-open topology is metrizable
using a countable compact exhaustion of $\HH$, so this sequential
compactness proves compactness.

Cauchy's formula gives locally uniform convergence of all derivatives.
The limiting derivative is nonzero, so reciprocals also converge locally
uniformly, as do their derivatives.  These observations give the asserted
joint continuity.  Compactness of $\KK\times E$ then gives the uniform
bounds.

For the explicit estimate, put $D=|x+i(y+1)|^2$ and $r=|m(x+iy)|$.
Then
\[
 1-r^2=\frac{4y}{D},\qquad
 |Q_F(x+iy)|=\frac{D}{4|f'(m(x+iy))|}.
\]
The lower derivative bound in \eqref{prelim:disk-bounds} gives
\[
 |Q_F(x+iy)|
 \le\frac{D(1+r)^3}{4(1-r)}
 =\frac{D^2(1+r)^4}{16y}
 \le\frac{D^2}{y}.
\]
When $|x|\le C_0$ and $y\le1$, this proves
\eqref{prelim:strip-bound}.

Finally $|m(z)|<1/7$ on $B(i,1/4)$, while
\eqref{prelim:quarter} gives
$B(0,1/16)\subset F(B(i,1/4))$.  Injectivity prevents any point outside
$B(i,1/4)$ from having an image of modulus less than $1/16$.  This proves
\eqref{prelim:zero-separation}.
\end{proof}

\subsection{Affine rescaling and dyadic growth}

Identify $z=x+iy\in\HH$ with the affine map $w\mapsto x+yw$.
Composition of these source maps is the product
\begin{equation}
 z\cdot w=x+yw,
 \qquad z^\#=\frac{-x+i}{y},
 \label{prelim:affine-product}
\end{equation}
with identity $i$ and inverse $z^\#$.  For $F\in\KK$ define
\begin{equation}
 T_zF(w)=\frac{F(z\cdot w)-F(z)}{yF'(z)}.
 \label{prelim:rerooting}
\end{equation}
It belongs to $\KK$.  On the space $C(\KK)$ of continuous
complex-valued functions, equipped with its supremum norm, define
\begin{equation}
 A_z\Phi(F)=|Q_F(z)|^2\Phi(T_zF),\qquad
 L=\frac12\bigl(A_{-1/2+i/2}+A_{1/2+i/2}\bigr).
 \label{prelim:operators}
\end{equation}
An operator is called positive here if it takes nonnegative real-valued
functions to nonnegative functions.  The operators in
\eqref{prelim:operators} are positive.  Their weights measure how the
inverse derivative changes when the source is rescaled.

\begin{lemma}[Affine identities and transfer growth]
\label{lem:transfer-growth}
For $z,w\in\HH$ and $F\in\KK$,
\begin{equation}
 Q_{T_zF}(w)=\frac{Q_F(z\cdot w)}{Q_F(z)},\qquad
 T_w(T_zF)=T_{z\cdot w}F,\qquad A_zA_w=A_{z\cdot w}.
 \label{prelim:composition}
\end{equation}
Each $A_z$ is bounded on $C(\KK)$, and for fixed $\Phi\in C(\KK)$
the map $z\mapsto A_z\Phi$ is continuous in the supremum norm.
For $n\ge0$, put $v_n=2^{-n}$ and
$u_{n,j}=-1+(2j+1)v_n$, $0\le j<2^n$.  Then
\begin{equation}
 L^n=v_n\sum_{j=0}^{2^n-1} A_{u_{n,j}+iv_n},\qquad
 \|L^n\|=\max_{F\in\KK}L^n\mathbf1(F).
 \label{prelim:dyadic-grid}
\end{equation}
The limit
\begin{equation}
 \rho=\lim_{n\to\infty}\|L^n\|^{1/n}
      =\inf_{n\ge1}\|L^n\|^{1/n}
 \label{prelim:rho}
\end{equation}
exists and satisfies $1\le\rho\le4$.  If $\lambda>\rho$, there is
$C_\lambda<\infty$ such that $\|L^n\|\le C_\lambda\lambda^n$ for
all $n\ge0$.
\end{lemma}

\begin{proof}
Differentiating \eqref{prelim:rerooting} proves the first identity in
\eqref{prelim:composition}.  Applying the normalization twice proves
the second.  In the operator product the weights telescope:
\[
 |Q_F(z)|^2|Q_{T_zF}(w)|^2=|Q_F(z\cdot w)|^2.
\]
This gives the third identity, with its displayed order of multiplication.
Joint continuity of $(z,F)\mapsto T_zF$ follows from locally uniform
convergence on affine images of compact source sets and the nonzero
normalizing denominator.  Together with Lemma~\ref{lem:compact-class},
this proves boundedness of $A_z$.  For fixed $\Phi$, joint continuity
on $\KK$ times a compact neighborhood of $z$ gives the asserted
supremum-norm continuity.

A word of length $n$ in the affine maps $-1/2+i/2$ and $1/2+i/2$
has height $2^{-n}$ and real part
$\sum_{l=1}^n\sigma_l2^{-l}$, where $\sigma_l\in\{-1,1\}$.
Writing $\sigma_l=2\epsilon_l-1$ and
$j=\sum_{l=1}^n\epsilon_l2^{n-l}$ gives real part
$-1+(2j+1)2^{-n}$.  Each $j$ occurs once and each word has coefficient
$2^{-n}$.  Expansion of $L^n$ proves the grid formula; at $n=0$ it
reduces to $A_i=I$.

For the positive sums defining $L^n$,
\[
 |L^n\Phi|\le L^n|\Phi|\le\|\Phi\|_\infty L^n\mathbf1.
\]
Testing with $\Phi=\mathbf1$ proves the norm identity.  At the map
$F(z)=z-i$ every summand in $L^n\mathbf1(F)$ equals one, so
$\|L^n\|\ge1$.  Since $|u_{n,j}|\le1$,
\eqref{prelim:strip-bound} gives
\[
 \|L^n\|\le625\,4^n.
\]
The logarithms of these positive norms are subadditive.  More explicitly,
for fixed $N\ge1$ write $n=qN+r$, $0\le r<N$, and use
$\|L^n\|\le\|L^N\|^q\|L^r\|$.  Taking the upper limit of the
$n$th roots and then the infimum over $N$ proves
\eqref{prelim:rho}; the displayed lower and upper bounds give
$1\le\rho\le4$.  If $\lambda>\rho$, the root limit bounds
$\|L^n\|$ by $\lambda^n$ for all sufficiently large $n$.  Increasing
a constant to cover the finitely many remaining $n$ proves the last claim.
\end{proof}

\subsection{From dyadic samples to disk integral means}

The next proposition explains why an upper bound on $\rho$ controls an
entire circle.  The comparison uses the same compact family in every
cell, so all its constants are independent of the chosen univalent map.

\begin{proposition}[Uniform disk comparison]
\label{prop:disk-transfer}
If $\rho\le2$, then for every $\varepsilon>0$ there is
$C_\varepsilon<\infty$ such that
\begin{equation}
 M_{-2}[f'](r)\le C_\varepsilon(1-r)^{-1-\varepsilon}
 \quad\bigl(f\in\SSS,\ 1/2\le r<1\bigr).
 \label{prelim:uniform-disk}
\end{equation}
The constant depends only on $\varepsilon$.
\end{proposition}

\begin{proof}
Fix $f\in\SSS$ and $1/2\le r<1$.  Put $\delta=1-r$ and choose
$n\ge1$ with $v=2^{-n}\le\delta<2v$.  On the fixed angular interval
$I=[3\pi/4,5\pi/4]$, write
\[
 m^{-1}(re^{i\theta})=x(\theta)+iy(\theta),\qquad
 D(\theta)=1-2r\cos\theta+r^2.
\]
The inverse formula gives
\begin{equation}
 x(\theta)=\frac{-2r\sin\theta}{D(\theta)},\qquad
 y(\theta)=\frac{1-r^2}{D(\theta)}.
 \label{prelim:arc-coordinates}
\end{equation}
On $I$ we have $\cos\theta\le0$ and hence $D\ge1+r^2\ge2r$;
therefore $|x|\le|\sin\theta|<1$.  Also
\[
 \frac\delta2\le y\le2\delta,\qquad
 \frac{dx}{d\theta}
 =\frac{4r^2-2r(1+r^2)\cos\theta}{D(\theta)^2}
 \ge\frac1{16}.
\]
The height bounds use $D\le(1+r)^2$, $D\ge1$, and $1+r\le2$.
For the derivative bound use $4r^2\ge1$ and $D^2\le16$.

The cells $[u_{n,j}-v,u_{n,j}+v]$ cover $[-1,1]$ with disjoint
interiors.  If $x(\theta)$ belongs to cell $j$, then
\[
 w=\frac{x(\theta)-u_{n,j}}v+i\frac{y(\theta)}v
 \in E:=\{u+i\eta:|u|\le1,\ 1/2\le\eta\le4\}.
\]
Let $B_E=\max_{G\in\KK,\,w\in E}|Q_G(w)|$, which is finite by
Lemma~\ref{lem:compact-class}.  The first identity in
\eqref{prelim:composition} gives
\begin{equation}
 |Q_F(x(\theta)+iy(\theta))|
 \le B_E|Q_F(u_{n,j}+iv)|,
 \qquad F=2i f\circ m.
 \label{prelim:cell-comparison}
\end{equation}
Thus the comparison covers every point of the varying-height arc in
each cell.

The exact derivative conversion in \eqref{prelim:normalization} is
\[
 |f'(m(z))|^{-2}
 =\left|\frac{m'(z)}{m'(i)}\right|^2|Q_F(z)|^2
 =\frac{16}{|z+i|^4}|Q_F(z)|^2
 \le16|Q_F(z)|^2.
\]
Since $x(\theta)$ is increasing and $d\theta\le16\,dx$,
\eqref{prelim:cell-comparison} and the cell lengths $2v$ give
\begin{align*}
 \int_I|f'(re^{i\theta})|^{-2}\,d\theta
 &\le256B_E^2\sum_j|Q_F(u_{n,j}+iv)|^2
       \bigl|[u_{n,j}-v,u_{n,j}+v]\cap x(I)\bigr|\\
 &\le512B_E^2 v\sum_j|Q_F(u_{n,j}+iv)|^2
 =512B_E^2 L^n\mathbf1(F).
\end{align*}
Here the bars around the intersection denote interval length.

For $\alpha\in\{1,i,-1,-i\}$, the map
$f_\alpha(z)=\alpha^{-1}f(\alpha z)$ belongs to $\SSS$ and has derivative
$f'(\alpha z)$.  The four rotated copies of $I$ cover the full circle,
with overlap only at endpoints.  Applying the preceding estimate to
these four maps therefore proves
\[
 M_{-2}[f'](r)\le C\|L^n\|
\]
with an absolute constant $C$.  If $\rho\le2$, apply
Lemma~\ref{lem:transfer-growth} with $\lambda=2^{1+\varepsilon}$.
It gives
\[
 M_{-2}[f'](r)\le C C_\lambda v^{-1-\varepsilon}
 \le C C_\lambda 2^{1+\varepsilon}(1-r)^{-1-\varepsilon}.
\]
The operator $L$ and the compact set $E$ were fixed before $f$ and $r$
were chosen, so this constant depends only on $\varepsilon$.
\end{proof}

We have reduced the uniform disk estimate to $\rho\le2$.  The next
sections assume $\rho>2$ and derive a contradiction from the probability
law generated by the same positive operator.
